% =========================================================
%  config/commands.tex
%  Perintah khusus dokumen RPS
%  Tujuan: mencegah definisi berulang yang tersebar di sections/.
% =========================================================

% --- Identitas mata kuliah -----------------------------
\newcommand{\NamaMK}{Pengantar Pemrograman Game}
\newcommand{\KodeMK}{MWCA009}
\newcommand{\NamaProdi}{Teknik Informatika}
\newcommand{\NamaFakultas}{Teknologi Informasi}
\newcommand{\NamaPT}{Universitas Bale Bandung}
\newcommand{\NamaDosen}{Cecep Suwanda, S.Si., M.Kom.}
\newcommand{\SemesterMK}{Ganjil}
\newcommand{\TahunAkademik}{2026/2027}

% --- Kode capaian pembelajaran -------------------------
\newcommand{\CPL}[1]{\textbf{CPL#1}}
\newcommand{\CPMK}[1]{\textbf{CPMK#1}}
\newcommand{\SUBCPMK}[1]{\textbf{Sub-CPMK#1}}

% --- Istilah teknis ------------------------------------
\newcommand{\code}[1]{\texttt{#1}}
\newcommand{\istilah}[1]{\textit{#1}}

% --- Penanda baris header tabel ------------------------
% Tanpa colortbl, penanda header cukup berupa penebalan garis pemisah
% (lihat \renewcommand{\toprule} pada config/styles.tex).
\newcommand{\headrow}{}

% --- Penanda status hubungan (matriks CPL--CPMK) -------
\newcommand{\Utama}{\textbf{Utama}}
\newcommand{\Pendukung}{Pendukung}
\newcommand{\TidakAda}{---}

% --- Opsi: mengaktifkan kembali latar baris tabel ------
% Apabila colortbl dapat dipakai pada instalasi LaTeX yang digunakan
% (yaitu array.sty sudah cukup baru), ganti \usepackage{xcolor} menjadi
% \usepackage[table]{xcolor} pada config/packages.tex, lalu ganti
% definisi di bawah ini:
%
%   \renewcommand{\headrow}{\rowcolor{ubHead}}
%   \renewcommand{\Utama}{\cellcolor{ubHead}\textbf{Utama}}
%   \renewcommand{\Pendukung}{\cellcolor{ubAlt}Pendukung}
